\begin{figure}[htbp]
\centering
\begin{tikzpicture}[font=\small,>=Stealth]
\fill[gray!7] (0,0) rectangle(2,2);
\draw[red!75!black,very thick,->] (0,0)--(0,2);
\draw[red!75!black,very thick,->] (2,0)--(2,2);
\draw[blue!75!black,very thick,->] (0,0)--(2,0);
\draw[blue!75!black,very thick,->] (2,2)--(0,2);
\node[below] at (1,0) {$x$};\node[left] at (0,1) {$y$};
\draw[->,thick] (2.7,1)--(4,1);
\begin{scope}[shift={(5.6,.2)},scale=.8]
% Shared bottle representation; the caller supplies scale and seam markers.
\fill[gray!8] (-.45,1.85) .. controls (-.5,1.1) and (-1,.9) .. (-1,.25)
 .. controls (-1,-.6) and (1,-.6) .. (1,.25)
 .. controls (1,.9) and (.45,1.1) .. (.45,1.85) --cycle;
\draw[thick] (-.45,1.85) .. controls (-.5,1.1) and (-1,.9) .. (-1,.25)
 .. controls (-1,-.6) and (1,-.6) .. (1,.25)
 .. controls (1,.9) and (.45,1.1) .. (.45,1.85);
\draw[thick] (-.45,1.85) .. controls (-.45,3.05) and (2.1,2.9) .. (2.1,1.5)
 .. controls (2.1,.9) and (.5,.8) .. (.15,.15);
\draw[thick] (.45,1.85) .. controls (.45,2.35) and (1.5,2.35) .. (1.5,1.5)
 .. controls (1.5,1.15) and (.25,1) .. (-.15,.15);

\draw[blue!75!black,very thick] (0,.12) ellipse(.2 and .1);
\node[below,blue!75!black] at (0,-.45) {glued seam};
\draw[orange!80!black,dotted,thick] (.85,.9) circle(.4);
\draw[->] (1.2,.9)--(2.35,.55);
\node[right,align=left] at (2.35,.55) {crossing only:\\no gluing};
\end{scope}
\end{tikzpicture}
\caption{Klein quotient: $(0,y)\sim(1,y)$ and $(x,0)\sim(1-x,1)$.
The solid blue seam represents the joined circles; the dotted callout
marks an incidental crossing of distinct sheets. The drawing in $\R^3$
is not an embedding. See Figure~\ref{fig:klein-gluing} for the deformation
and Figure~\ref{fig:klein-corner} for the local disk at the corner class.}
\label{fig:quotient-klein}
\end{figure}
